\subsection{Cubic derivative equations also prescribe a unique singular jet}
\label{sec:hermite-johnson-cubic}

A nonzero cubic has at most one distinct repeated root.  When such a
root persists with the challenge, it is rational in the coefficients.
This extends the preceding argument without selecting among several
possible derivative branches.

\begin{corollary}[Cubic derivative identities]
\label{cor:hermite-johnson-cubic}
Retain the domain, received line, degree bound, and definition of bad
pairs from Theorem~\ref{thm:hermite-johnson-quadratic}.  Let
$Q(X,z,u,v)\ne0$ have derivative degree at most three, total $(u,v)$
degree at most $B_0$, and challenge degree at most $H_0$, with unrestricted
$X$-degree.  Put $h=B_0+H_0$ and
\[
 F_x(z,v)=Q(x,z,f_x+zg_x,v)\in F[z,v].
\]
Define $S_{\rm ord}$ by $F_x=0$ identically.  Define $S_{\rm jet}$
to consist of those other coordinates where the nonzero polynomial
$F_x\in F(z)[v]$ has a repeated root over an algebraic closure of
$F(z)$.  Write $u=|S_{\rm ord}|$ and $s=|S_{\rm jet}|$.

Suppose $D/n\ge\rho_{\min}>0$ and
\[
 A-D\ge\eta n,\qquad
 A-u\ge\sqrt{Ds/2}+\varepsilon n
\]
for fixed positive $\eta,\varepsilon$.  In characteristic zero, or
characteristic $p>\max\{D,3,B\}$ for the bounded interpolation degree
$B$ below, there are
$O_{\rho_{\min},\eta,\varepsilon,B_0,H_0}(n)$ bad labels admitting
an actual solution $Q(X,z,P,P')=0$.

More precisely, retain $L,e,b,M,T,B,V,c_M$ from
\eqref{eq:hj-interpolation-parameters}, with $V>sc_M$, and replace
\eqref{eq:hj-challenge-degree} by
\begin{equation}
 H=\left\lfloor\frac{sc_MB(2h+1)}{V-sc_M}\right\rfloor.
 \label{eq:hj-cubic-challenge-degree}
\end{equation}
Then the finite bound~\eqref{eq:hj-finite-bound} holds with its
singular-incidence term $2dn/(e+1)$ replaced by $4hn/(e+1)$,
and with its algebraic terms evaluated at these new $B,H,T$.
The regular-incidence term still uses the original $B_0,H_0$.
\end{corollary}

\begin{proof}
Every coefficient of $F_x$ has challenge degree at most $h$.  We show
that outside at most $4h$ exceptional labels at each nonordinary
coordinate, every singular match belongs to $S_{\rm jet}$ and has
one prescribed rational derivative, of numerator and denominator
degrees at most $2h$.  All degree classifications below are generic
over $F(z)$; the exceptional sets explicitly include degree drops.

First suppose
\[
 F_x=a v^3+b v^2+c v+d,\qquad a\ne0.
\]
Set
\begin{align*}
 \Delta&=b^2c^2-4ac^3-4b^3d-27a^2d^2+18abcd,\\
 \Delta_0&=b^2-3ac,\qquad V_0=9ad-bc.
\end{align*}
Their challenge degrees are at most $4h,2h,2h$, respectively.
If $\Delta\ne0$ as a polynomial in $z$, every singular specialization
lies among its at most $4h$ roots.  This includes leading-coefficient
zeros: at $a(z)=0$ the expression becomes
$b(z)^2(c(z)^2-4b(z)d(z))$, which vanishes at a repeated quadratic
root and at any further degree drop admitting a singular root.

If $\Delta=0$ identically and $\Delta_0\ne0$, exclude the roots of
$a\Delta_0$, at most $3h$ labels.  At every remaining label the unique
repeated root is
\begin{equation}
 r_x(z)=\frac{V_0(z)}{2\Delta_0(z)}.
 \label{eq:hj-cubic-repeated-root}
\end{equation}
Indeed, writing the cubic over an algebraic closure as
$a(v-r)^2(v-t)$ with $r\ne t$ gives
$\Delta_0=a^2(r-t)^2$ and $V_0=2a^2r(r-t)^2$.
If both $\Delta$ and $\Delta_0$ vanish identically, exclude the roots
of $a$.  The remaining polynomial is
$a(v+b/(3a))^3$, and its unique repeated root is $-b/(3a)$.
This follows from $c=b^2/(3a)$ and $d=b^3/(27a^2)$, using
characteristic different from two and three.  These are exactly the
persistent repeated-root cases of generic cubic degree; all their
prescriptions have rational height at most $2h$.

If the generic degree is two, write $F_x=bv^2+cv+d$, with $b\ne0$.
A nonzero discriminant $c^2-4bd$ confines every singular specialization,
including degree drops, to at most $2h$ labels.  An identically zero
discriminant prescribes the unique repeated root $-c/(2b)$ outside
the at most $h$ roots of $b$.
For generic degree one, singular agreement requires the nonzero linear
coefficient to vanish.  For a nonzero constant polynomial in $v$, an
actual agreement requires that constant coefficient to vanish.  Each
case contributes at most $h$ exceptional labels and no persistent
jet core.  The identically zero case is precisely $S_{\rm ord}$.

Thus the exceptional coordinate--label incidence budget is $4hn$,
uniformly over all candidates.  After the regular-incidence deletion
and deletion of at most $4hn/(e+1)$ further labels, every remaining
bad witness has at least $A-u-(L-1)-e$ matches in $S_{\rm jet}$ with
$P'(x)=r_x(z)$ defined.  Write $r_x=N_x/E_x$ with both degrees at
most $2h$.  In the local Hermite expansion multiply by $E_x^B$.
Each cleared constraint now has challenge degree at most
$H+B(2h+1)$, so~\eqref{eq:hj-cubic-challenge-degree} gives
\[
 (H+1)V>sc_M\bigl(H+B(2h+1)+1\bigr).
\]
The multiplicity argument, specialization-safe algebraic-root lemma,
and original full-support incidence count are unchanged.  The same
fixed-margin choices of $L,e,M$ make $B,H$ bounded and give the
claimed linear bound.
\end{proof}

The assertion is not that every cubic has a repeated derivative root:
generically squarefree specializations contribute only the bounded
exceptional incidences.  It is the \emph{uniqueness and rationality}
of a persistent repeated root that permits this extension.  Degree
four is the first degree allowing two distinct repeated derivative
roots, as in $(v-r_1)^2(v-r_2)^2$.  The corollary leaves the ordinary
core, the agreement margin, and containment in one common bounded-degree
identity as explicit hypotheses; it is not an unconditional obstruction
to all proximity-gap counterexamples.
